\documentclass{exptracker}
% Add the `final` option when the report is ready to share:
%   \documentclass[final]{exptracker}
% It hides the grey guidance notes and warns in the log about any
% \placeholder left unfilled.

\title{Experiment title}
\author{First Author \and Second Author}
\date{\today}
\version{0.1}
% Draft, Proposed, In progress, Complete or Abandoned.
\status{Draft}
% The report's Overleaf project URL. Leave it empty until the report is on
% Overleaf; /upload-overleaf-experiment-report fills it in.
\overleafproject{}
% One \relatedlink per link (GitHub issue or PR, W&B run, Azure ML job, eval
% outputs). Paste URLs as they are: & % # _ need no escaping.
\relatedlink{GitHub issue}{https://github.com/}
\relatedlink{W\&B project}{https://wandb.ai/}

\begin{document}
\maketitle

\section{Executive summary}
\guidance{The key findings, for a reader who reads nothing else. Point each one at
the result it summarises with \texttt{\textbackslash ref}.}

\begin{itemize}
  \item \placeholder{A key finding in a sentence or two, with its headline number,
    e.g.\ ``doing X improves Y by Z (\ref{find:main})''.}
\end{itemize}

\section{Problem statement}
\guidance{Briefly describe the problem and why it matters.}

\placeholder{A paragraph or two: what is wrong or missing, who it affects, and
why it is worth an experiment.}

\section{Hypotheses}
\guidance{What you expect to happen. Give each hypothesis a
\texttt{\textbackslash label} so later sections can cite it with
\texttt{\textbackslash ref}, which prints its ID.}

\begin{hypotheses}
  \item\label{hyp:main} \placeholder{A testable statement, e.g.\ ``doing X
    improves Y by at least Z on W''.}
  \item\label{hyp:secondary} \placeholder{Another hypothesis, or delete this
    item.}
\end{hypotheses}

\section{Assumptions}
\guidance{Key assumptions underlying the hypotheses.}

\begin{assumptions}
  \item \placeholder{Something taken as given that, if false, would undermine
    \ref{hyp:main}.}
\end{assumptions}

\section{Proposed approach and model choices}
\guidance{What you plan to do and why, and what else you considered.}

\begin{approaches}
  \item \placeholder{What will be done.}
    \lead{Rationale} \placeholder{Why this choice.}
    \lead{Alternatives considered} \placeholder{What else was on the table and
    why it was not chosen.}
\end{approaches}

\section{Supporting evidence}
\guidance{Why the hypotheses are reasonable: prior results, literature, pilot
runs.}

\begin{evidence}
  \item \placeholder{A prior result or paper that supports \ref{hyp:main},
    e.g.\ \citet{placeholder2026}.}
\end{evidence}

\section{Evaluation}
\guidance{How success will be measured: metrics, baselines, data and
thresholds.}

\begin{criteria}
  \item \placeholder{Metric, baseline and threshold that would support
    \ref{hyp:main}.}
\end{criteria}

\section{Risks and limitations}
\guidance{Known concerns or constraints.}

\begin{risks}
  \item \placeholder{A risk or limitation, and how it will be mitigated.}
\end{risks}

\section{Expected impact}
\guidance{What happens if each hypothesis is supported or rejected?}

\begin{impacts}
  \item \placeholder{The decision this informs.}
    \lead{If \ref{hyp:main} is supported} \placeholder{What changes.}
    \lead{If \ref{hyp:main} is rejected} \placeholder{What changes.}
\end{impacts}

\section{Results}
\guidance{What the experiments showed, whether each hypothesis held, and the
follow-ups. Verdicts: Supported, Partly supported, Inconclusive, Rejected or
Pending.}

\begin{findings}
  \item\label{find:main} \verdict{Pending} \ref{hyp:main}: \placeholder{The result, with numbers
    and a pointer to the figure or table that shows it.}
    \lead{Follow-up} \placeholder{What to do next.}
\end{findings}

% Figures go where they are discussed, e.g.
% \begin{figure}[ht]
%   \centering
%   \includegraphics[width=0.8\linewidth]{figures/crps.pdf}
%   \caption{What the figure shows.}\label{fig:crps}
% \end{figure}

% Delete these two lines if the report cites nothing.
\bibliographystyle{plainnat}
\bibliography{references}

\end{document}
